
We asked: does when you measure GPU memory matter as much as how many times you measure? Our results provide directional evidence---post-optimizer timing achieves 52\% error reduction compared to pre-optimizer sampling across 4 CNN configurations, demonstrating that sampling timing is as critical as iteration count for accurate memory profiling.

This work reframes lightweight memory profiling from purely an iteration-count problem to a timing problem. By identifying that optimizer workspace allocations occur during \texttt{optimizer.step()} rather than during backward propagation, we designed a 3-iteration post-optimizer sampling protocol that captures the precise moment when Adam's momentum buffers are allocated in GPU memory. Across ResNet-18 and ResNet-34 architectures with Adam and SGD optimizers, our protocol achieves 2.6\% median relative error---all tested configurations remain below 7\%, well within acceptable bounds for pre-training OOM prediction in CNN training scenarios.

The optimizer-specific accuracy finding---Adam achieving $17\times$ lower error than SGD despite smaller workspace---reveals that memory profiling is not optimizer-agnostic. This suggests optimizer-specific timing profiles require investigation. Understanding SGD's allocation timing characteristics is critical to validate whether post-optimizer sampling is a general mechanism or requires optimizer-specific calibration.

Beyond the quantitative improvements, our fast validation methodology demonstrates practical research value. Validating core mechanisms with 4 configurations in approximately 4 seconds---rather than committing to 48 configurations over GPU-hours---enabled rapid hypothesis iteration during development.

Just as measuring highway traffic at 3 PM misses rush hour at 5 PM, sampling memory before \texttt{optimizer.step()} misses the allocation peak that determines OOM failures. By capturing the right moment---immediately after optimizer workspace allocation completes---we enable accurate lightweight memory profiling for CNN architectures. A researcher training ResNet-34 with batch size 64 can now predict in 3 iterations whether full training will encounter OOM, saving GPU-hours otherwise wasted on failed experiments.

Three immediate extensions would strengthen and broaden our findings. First, transformer validation across architectures with variable-length sequence inputs and stratified sampling would test whether post-optimizer timing generalizes beyond CNNs. Second, investigating SGD's timing characteristics at microsecond granularity could reveal why SGD achieves $17\times$ higher error than Adam and inform design of SGD-optimized protocols. Third, validating with real datasets including DataLoader overhead would confirm that our synthetic data findings transfer to production training scenarios.

The post-optimizer sampling principle---measure at the moment when critical allocations complete---extends beyond memory profiling. Any resource prediction system must consider not just what to measure but when to measure it relative to system state transitions. This work demonstrates that temporal precision in measurement timing can yield substantial accuracy improvements, a lesson applicable to profiling systems beyond GPU memory prediction.

